%% kapa: exam=2 round=25 year=2014 period=2 subject=appraisal_theory subject_ko=감정평가이론 seq=2 number=2 points=35
%% source: Q-Net 감정평가사 기출문제 (cst00309) artlSeq=5109209 file=extracted/1008/5109209/2014년 제25회 감정평가사 2차시험 기출문제/2교시(감정평가이론).hwp
%% conversion: kordoc 4.24.51 -> tools/gen2.py (auto); review=proofread-sonnet; flags=none
\begin{longquestion}{【문제 2】}
근린형 쇼핑센터 내 구분점포(「집합건물의 소유 및 관리에 관한 법률」에 의한 상가건물의 구분소유 부분)의 시장가치를 감정평가하려 한다. 인근에 경쟁적인 초대형 쇼핑센터가 입지하여, 대상점포가 소재한 근린형 쇼핑센터의 고객흡인력이 급격히 감소하고 상권이 위축되어 구분점포 거래가 희소하게 된 시장동향을 고려하여 다음 물음에 답하시오. (35점) \par
1) 대상 구분점포의 감정평가에 거래사례비교법을 적용할 경우 감정평가방법의 개요, 적용상 한계 및 수집된 거래사례의 거래조건보정(Transactional adjustments)에 대하여 설명하고, 그 밖에 적용 가능한 다른 감정평가방법의 개요 및 적용시 유의할 사항에 대하여 설명하시오. (25점) \par
2) 적용된 각 감정평가방법에 의한 시산가액 간에 괴리가 발생되었을 경우 시산가액 조정의 의미, 기준 및 재검토할 사항에 대하여 설명하시오. (10점) \par
\end{longquestion}
